
\section{Advanced Activities}

The following activities are intended for strong master's or PhD students.
Students may use GPOPS-II, Dymos, or an equivalent optimal-control and
direct-collocation package where specified.

\subsection*{Problem 1: Perfect-Information OLOC versus Causal Feedback}

Consider
\begin{equation}
\dot{x}(t)=-0.8x(t)+u(t)+w(t),
\qquad
x(0)=1,
\qquad
0\leq t\leq 6,
\end{equation}
with
\begin{equation}
-3\leq u(t)\leq 3.
\end{equation}

The nominal disturbance is
\begin{equation}
w_0(t)=0.4\sin(1.5t)+0.2\sin(4t).
\end{equation}

Minimize
\begin{equation}
J=
\frac{1}{2}x(6)^2
+
\frac{1}{2}
\int_0^6
\left[
x(t)^2+0.05u(t)^2
\right]dt.
\end{equation}

\begin{enumerate}
    \item Write the Hamiltonian and derive the costate equation.

    \item Show that the unconstrained optimal control is
    \begin{equation}
    u^*(t)=-\frac{\lambda(t)}{0.05},
    \end{equation}
    and write the corresponding saturated optimal-control law.

    \item State the terminal condition for the costate.

    \item Solve the two-point boundary-value problem using GPOPS-II, Dymos,
    or a shooting method.

    \item Design a finite-horizon causal LQR controller
    \begin{equation}
    u(t)=-K(t)x(t)
    \end{equation}
    by solving the Riccati equation backward in time.

    \item Apply both controllers to the perturbed disturbance
    \begin{equation}
    w(t)=w_0(t)+0.15\sin(2.3t+\phi),
    \end{equation}
    for
    \begin{equation}
    \phi\in
    \left\{
    0,\frac{\pi}{2},\pi,\frac{3\pi}{2}
    \right\}.
    \end{equation}

    \item Compare the nominal and perturbed costs. Quantify the performance
    advantage caused by perfect future disturbance information.
\end{enumerate}


\subsection*{Problem 2: Open-Loop versus Feedback Control Co-Design}

Consider
\begin{equation}
\dot{x}_1=x_2,
\end{equation}
\begin{equation}
\dot{x}_2=-kx_1-cx_2+u+d(t),
\end{equation}
where the mass is normalized to one.

The plant variables are
\begin{equation}
0.5\leq k\leq 8,
\qquad
0.1\leq c\leq 3,
\end{equation}
and the actuator-capacity variable is
\begin{equation}
0.5\leq F_{\max}\leq 5.
\end{equation}

Use the feedback law
\begin{equation}
u(t)=
\operatorname{sat}
\left(
-K_px_1(t)-K_dx_2(t),
F_{\max}
\right),
\end{equation}
with
\begin{equation}
0\leq K_p\leq 10,
\qquad
0\leq K_d\leq 8.
\end{equation}

Minimize
\begin{align}
J=&
\frac{1}{N_s}
\sum_{s=1}^{N_s}
\int_0^8
\left[
x_{1,s}^2
+0.1x_{2,s}^2
+0.02u_s^2
\right]dt
\nonumber\\
&+
0.01k^2
+
0.02c^2
+
0.04F_{\max}^2.
\end{align}

Use the scenarios
\begin{equation}
\mathbf{x}_0^{(1)}=
\begin{bmatrix}
1\\
0
\end{bmatrix},
\qquad
\mathbf{x}_0^{(2)}=
\begin{bmatrix}
0.5\\
1
\end{bmatrix},
\qquad
\mathbf{x}_0^{(3)}=
\begin{bmatrix}
-1\\
0.5
\end{bmatrix},
\end{equation}
with
\begin{equation}
d_1(t)=0,
\qquad
d_2(t)=0.4\sin(2t),
\qquad
d_3(t)=0.3\sin(3t).
\end{equation}

\begin{enumerate}
    \item Derive the unsaturated closed-loop dynamics.

    \item Derive the conditions on $k$, $c$, $K_p$, and $K_d$ for
    asymptotic stability.

    \item Add the passive-safety constraint
    \begin{equation}
    \zeta_{\mathrm{passive}}
    =
    \frac{c}{2\sqrt{k}}
    \geq 0.15.
    \end{equation}

    \item Optimize
    \begin{equation}
    k,\quad c,\quad F_{\max},\quad K_p,\quad K_d
    \end{equation}
    simultaneously.

    \item For the optimized plant, solve a separate OLOC problem in
    GPOPS-II or Dymos for each disturbance scenario.

    \item Compute the relative implementation gap
    \begin{equation}
    G=
    \frac{J_{\mathrm{feedback}}-J_{\mathrm{OLOC}}}
    {J_{\mathrm{OLOC}}}.
    \end{equation}

    \item Explain why the plant optimized with unrestricted OLOC may differ
    from the plant optimized with the parameterized feedback controller.
\end{enumerate}


\subsection*{Problem 3: Disturbance Preview and Information Availability}

Consider the discrete-time system
\begin{equation}
\mathbf{x}_{k+1}
=
A\mathbf{x}_k
+
Bu_k
+
Ew_k,
\end{equation}
where
\begin{equation}
A=
\begin{bmatrix}
1 & h\\
0 & 1
\end{bmatrix},
\qquad
B=E=
\begin{bmatrix}
h^2/2\\
h
\end{bmatrix},
\qquad
h=0.1.
\end{equation}

Use the horizon
\begin{equation}
N=60
\end{equation}
and the objective
\begin{equation}
J=
\mathbf{x}_N^TQ_f\mathbf{x}_N
+
\sum_{k=0}^{N-1}
\left(
\mathbf{x}_k^TQ\mathbf{x}_k
+
Ru_k^2
\right),
\end{equation}
with
\begin{equation}
Q=
\begin{bmatrix}
10&0\\
0&1
\end{bmatrix},
\qquad
Q_f=10Q,
\qquad
R=0.1.
\end{equation}

The disturbance is
\begin{equation}
w_k=
0.8\sin(0.25k)
+
0.3\sin(0.9k).
\end{equation}

Assume that at time $k$, the controller knows only the next $L$
disturbance values.

\begin{enumerate}
    \item Assume the value function has the affine-quadratic form
    \begin{equation}
    V_k(\mathbf{x})
    =
    \mathbf{x}^TP_k\mathbf{x}
    +
    2\mathbf{s}_k^T\mathbf{x}
    +
    c_k.
    \end{equation}
    Derive the dynamic-programming recursion for $P_k$, $\mathbf{s}_k$,
    and $c_k$.

    \item Show that the optimal control has the form
    \begin{equation}
    u_k
    =
    -K_k\mathbf{x}_k
    -
    u_{k,\mathrm{ff}},
    \end{equation}
    where the feedforward term depends on the available disturbance preview.

    \item Implement controllers for
    \begin{equation}
    L=0,\ 2,\ 5,\ 10,\ 20,\ 60.
    \end{equation}

    \item Plot the optimal cost as a function of preview length.

    \item Compute the marginal value of preview:
    \begin{equation}
    \Delta J_L=J_L-J_{L+1}.
    \end{equation}

    \item Determine the preview length beyond which additional information
    provides less than a $1\%$ improvement.

    \item Explain how the optimal physical design of a suspension or energy
    converter could change as $L$ changes.
\end{enumerate}


\subsection*{Problem 4: Sensor and Output-Feedback Co-Design}

Consider the stochastic plant
\begin{equation}
\dot{\mathbf{x}}
=
A(p)\mathbf{x}
+
Bu
+
\mathbf{w},
\end{equation}
where
\begin{equation}
A(p)=
\begin{bmatrix}
0&1\\
-p&-0.4
\end{bmatrix},
\qquad
B=
\begin{bmatrix}
0\\
1
\end{bmatrix},
\end{equation}
and
\begin{equation}
1\leq p\leq4
\end{equation}
is a plant-design variable.

The process-noise covariance is
\begin{equation}
W=
\begin{bmatrix}
0&0\\
0&0.04
\end{bmatrix}.
\end{equation}

Three sensor architectures are available:
\begin{equation}
C_1=
\begin{bmatrix}
1&0
\end{bmatrix},
\qquad
V_1=0.01,
\qquad
C_{\mathrm{sensor},1}=0.20,
\end{equation}
\begin{equation}
C_2=
\begin{bmatrix}
0&1
\end{bmatrix},
\qquad
V_2=0.04,
\qquad
C_{\mathrm{sensor},2}=0.10,
\end{equation}
and
\begin{equation}
C_3=I,
\qquad
V_3=
\begin{bmatrix}
0.01&0\\
0&0.04
\end{bmatrix},
\qquad
C_{\mathrm{sensor},3}=0.35.
\end{equation}

Use the control-performance weights
\begin{equation}
Q=
\begin{bmatrix}
10&0\\
0&1
\end{bmatrix},
\qquad
R=0.2.
\end{equation}

\begin{enumerate}
    \item Compute the observability matrix for each sensor architecture and
    determine whether the system is observable for all feasible $p$.

    \item For each value of $p$, solve the LQR algebraic Riccati equation and
    compute the state-feedback gain $K(p)$.

    \item For each sensor architecture, solve the Kalman-filter Riccati equation
    \begin{equation}
    A\Sigma+\Sigma A^T
    -
    \Sigma C^TV^{-1}C\Sigma
    +
    W=0
    \end{equation}
    and compute
    \begin{equation}
    L=\Sigma C^TV^{-1}.
    \end{equation}

    \item Construct the LQG controller
    \begin{equation}
    u=-K\hat{\mathbf{x}}.
    \end{equation}

    \item For each sensor architecture, optimize $p$ using
    \begin{equation}
    J_{\mathrm{total}}
    =
    \mathbb{E}
    \left[
    \int_0^{20}
    \left(
    \mathbf{x}^TQ\mathbf{x}
    +
    Ru^2
    \right)dt
    \right]
    +
    0.03p^2
    +
    C_{\mathrm{sensor}}.
    \end{equation}

    \item Estimate the expected cost using at least 200 Monte Carlo
    noise realizations.

    \item Identify the optimal plant--sensor--controller combination.

    \item Explain why the separation principle does not imply that sensor
    selection and plant design can be performed independently.
\end{enumerate}


\subsection*{Problem 5: Nested Plant--MPC Co-Design}

Consider the sampled system
\begin{equation}
\mathbf{x}_{k+1}
=
A(k_s)\mathbf{x}_k
+
Bu_k
+
Ew_k,
\end{equation}
where
\begin{equation}
A(k_s)=
\begin{bmatrix}
1&h\\
-k_sh&1-0.4h
\end{bmatrix},
\end{equation}
\begin{equation}
B=E=
\begin{bmatrix}
0\\
h
\end{bmatrix},
\qquad
h=0.1.
\end{equation}

The plant variables are
\begin{equation}
0.5\leq k_s\leq5,
\qquad
0.5\leq F_{\max}\leq3.
\end{equation}

At every time step, use MPC to solve
\begin{equation}
\min_{\{u_i\}}
\sum_{i=0}^{N_p-1}
\left(
\mathbf{x}_i^TQ\mathbf{x}_i
+
Ru_i^2
\right)
+
\mathbf{x}_{N_p}^TP\mathbf{x}_{N_p},
\end{equation}
subject to
\begin{equation}
\mathbf{x}_{i+1}=A(k_s)\mathbf{x}_i+Bu_i,
\end{equation}
\begin{equation}
|u_i|\leq F_{\max},
\qquad
|x_{1,i}|\leq1.
\end{equation}

Use
\begin{equation}
Q=
\begin{bmatrix}
10&0\\
0&1
\end{bmatrix},
\qquad
R=0.05.
\end{equation}

\begin{enumerate}
    \item Derive the condensed prediction equation
    \begin{equation}
    \mathbf{X}
    =
    \mathcal{A}\mathbf{x}_k
    +
    \mathcal{B}\mathbf{U}.
    \end{equation}

    \item Write the MPC problem as a quadratic program
    \begin{equation}
    \min_{\mathbf{U}}
    \frac{1}{2}\mathbf{U}^TH\mathbf{U}
    +
    \mathbf{f}^T\mathbf{U}.
    \end{equation}

    \item Implement a nested CCD architecture in which the outer optimization
    chooses $k_s$ and $F_{\max}$, while the inner problem performs
    closed-loop MPC simulation.

    \item Use prediction horizons
    \begin{equation}
    N_p=5,\ 10,\ 20,\ 40.
    \end{equation}

    \item Minimize
    \begin{equation}
    J_{\mathrm{CCD}}
    =
    \sum_{k=0}^{100}
    \left(
    \mathbf{x}_k^TQ\mathbf{x}_k
    +
    Ru_k^2
    \right)
    +
    0.02k_s^2
    +
    0.05F_{\max}^2.
    \end{equation}

    \item Compare the optimal plant for each prediction horizon.

    \item Compare the MPC result with a full-horizon OLOC solution obtained
    using GPOPS-II or Dymos.

    \item Introduce a $20\%$ stiffness-model error and determine whether
    recursive feasibility is maintained.

    \item Explain why prediction horizon should be regarded as a
    control-design variable in CCD.
\end{enumerate}


\subsection*{Problem 6: Converting an OLOC Trajectory into TVLQR Feedback}

Consider a torque-controlled pendulum. Let $\theta=0$ denote the upright
position and $\theta=\pi$ the downward position.

The dynamics are
\begin{equation}
\dot{\theta}=\omega,
\end{equation}
\begin{equation}
\dot{\omega}
=
\frac{g}{l}\sin\theta
-
\frac{b}{ml^2}\omega
+
\frac{1}{ml^2}u.
\end{equation}

Use
\begin{equation}
m=1\ {\rm kg},
\qquad
l=1\ {\rm m},
\qquad
b=0.05\ {\rm N\,m\,s},
\end{equation}
\begin{equation}
g=9.81\ {\rm m/s^2},
\qquad
|u(t)|\leq8\ {\rm N\,m}.
\end{equation}

The boundary conditions are
\begin{equation}
\mathbf{x}(0)=
\begin{bmatrix}
\pi\\
0
\end{bmatrix},
\qquad
\mathbf{x}(4)=
\begin{bmatrix}
0\\
0
\end{bmatrix}.
\end{equation}

Minimize
\begin{equation}
J=
\int_0^4
\left(
0.001\omega(t)^2
+
0.01u(t)^2
\right)dt.
\end{equation}

\begin{enumerate}
    \item Solve the swing-up OLOC problem using GPOPS-II or Dymos.

    \item Denote the optimal trajectory by
    \begin{equation}
    \mathbf{x}^*(t),
    \qquad
    u^*(t).
    \end{equation}
    Linearize the dynamics along the optimal trajectory and show that
    \begin{equation}
    A(t)=
    \begin{bmatrix}
    0&1\\
    \dfrac{g}{l}\cos\theta^*(t)&-\dfrac{b}{ml^2}
    \end{bmatrix},
    \qquad
    B(t)=
    \begin{bmatrix}
    0\\
    \dfrac{1}{ml^2}
    \end{bmatrix}.
    \end{equation}

    \item Solve the time-varying Riccati equation
    \begin{equation}
    -\dot{S}
    =
    A^TS+SA
    -
    SBR^{-1}B^TS
    +
    Q,
    \end{equation}
    with terminal condition
    \begin{equation}
    S(4)=Q_f.
    \end{equation}

    \item Implement the trajectory-tracking controller
    \begin{equation}
    u(t)
    =
    \operatorname{sat}
    \left[
    u^*(t)
    -
    R^{-1}B(t)^TS(t)
    \left(
    \mathbf{x}(t)-\mathbf{x}^*(t)
    \right)
    \right].
    \end{equation}

    \item Compare pure open-loop implementation and TVLQR feedback for
    the perturbed initial conditions
    \begin{equation}
    \theta(0)=\pi+\Delta\theta,
    \qquad
    \Delta\theta\in
    \{-0.15,-0.10,0.10,0.15\}.
    \end{equation}

    \item Repeat with
    \begin{equation}
    10\% \text{ error in }m,
    \qquad
    10\% \text{ error in }l,
    \end{equation}
    and additive measurement noise.

    \item Report
    \begin{equation}
    \|\mathbf{x}(4)\|_2,
    \qquad
    \max_t|u(t)|,
    \qquad
    J.
    \end{equation}

    \item Determine the largest initial perturbation for which the TVLQR
    controller successfully reaches the upright neighborhood
    \begin{equation}
    |\theta(4)|\leq0.05,
    \qquad
    |\omega(4)|\leq0.1.
    \end{equation}

    \item Explain why the OLOC trajectory remains valuable even though it is
    not robust enough for direct implementation.
\end{enumerate}
